\backmatter

\chapter{Index of Definitions}
The locators below retain the original chapter and section numbers and link to those sections in this edition.
\begingroup
\renewcommand{\arraystretch}{0.9}
\begin{longtable}{@{}p{0.75\textwidth} l@{}}
    acyclic complex & \bookref{sec:1.5}{I.5} \\
    bounded complex & \bookref{sec:1.2}{I.2} \\
    Cartan-Eilenberg resolution & \bookref{sec:1.7}{I.7} \\
    catenary & \bookref{sec:5.7}{V.7} \\
    \v{C}ech complex & \bookref{sec:3.3}{III.3} \\
    codimension & \bookref{sec:4.1}{IV.1} \\
    codimension function & \bookref{sec:5.7}{V.7} \\
    co-finite type module & \bookref{sec:5.5}{V.5} \\
    Cohen-Macaulay complex & \bookref{sec:4.3}{IV.3} \\
    Cohen-Macaulay morphism & \bookref{sec:5.9}{V.9} \\
    Cohen-Macaulay-$1$ ring & \bookref{sec:5.9}{V.9} \\
    Cohen-Macaulay sheaf & \bookref{sec:4.2}{IV.2} \\
    cohomological functor & \bookref{sec:1.1}{I.1} \\
    complex & \bookref{sec:1.2}{I.2} \\
    Cousin complex & \bookref{sec:4.3}{IV.3} \\
    Cousin complex of $\mathcal{F}$ & \bookref{sec:4.2}{IV.2} \\
    $\partial$-functor & \bookref{sec:1.1}{I.1} \\
    depth & \bookref{sec:4.2}{IV.2} \\
    derived category & \bookref{sec:1.4}{I.4} \\
    derived functor & \bookref{sec:1.5}{I.5} \\
    duality morphism & \bookref{sec:3.5}{III.5} \\
    dualizing complex & \bookref{sec:5.2}{V.2} \\
    dualizing functor & \bookref{sec:5.5}{V.5} \\
    embeddable morphism & \bookref{sec:3.8}{III.8} \\
    $F$-acyclic & \bookref{sec:1.5}{I.5} \\
    family of supports & \bookref{sec:4.1}{IV.1} \\
    filtration associated to a dualizing complex & \bookref{sec:5.7}{V.7} \\
    finite cohomological dimension & \bookref{sec:1.5}{I.5} \\
    finite injective dimension & \bookref{sec:1.7}{I.7} \\
    finite Tor-dimension & \bookref{sec:2.4}{II.4} \\
    Gorenstein complex & \bookref{sec:4.3}{IV.3} \\
    Gorenstein morphism & \bookref{sec:3.1}{III.1}, \bookref{sec:5.9}{V.9} \\
    Gorenstein prescheme & \bookref{sec:5.9}{V.9} \\
    Gorenstein ring & \bookref{sec:5.9}{V.9} \\
    homotopy & \bookref{sec:1.2}{I.2} \\
    hyperext & \bookref{sec:1.6}{I.6} \\
    Koszul complex & \bookref{sec:3.7}{III.7} \\
    Krull dimension & \bookref{sec:5.2}{V.2} \\
    local complete intersection & \bookref{sec:3.1}{III.1} \\
    local hyperext & \bookref{sec:2.3}{II.3} \\
    local hyper Tor & \bookref{sec:2.4}{II.4} \\
    localization of a category & \bookref{sec:1.3}{I.3} \\
    localizing subcategory & \bookref{sec:1.5}{I.5} \\
    locally noetherian category & \bookref{sec:2.7}{II.7} \\
    local parameter & \bookref{sec:7.1}{VII.1} \\
    mapping cone & \bookref{sec:1.2}{I.2} \\
    multiplicative system & \bookref{sec:1.3}{I.3} \\
    noetherian object & \bookref{sec:2.7}{II.7} \\
    normalized dualizing complex & \bookref{sec:5.6}{V.6} \\
    pointwise bounded below & \bookref{sec:5.8}{V.8} \\
    pointwise dualizing complex & \bookref{sec:5.8}{V.8} \\
    pointwise finite injective dimension & \bookref{sec:5.8}{V.8} \\
    projectively embeddable & \bookref{sec:3.10}{III.10} \\
    rational point & \bookref{sec:7.1}{VII.1} \\
    reflexive complex & \bookref{sec:5.2}{V.2} \\
    regular prescheme & \bookref{sec:5.2}{V.2} \\
    relative differentials & \bookref{sec:3.1}{III.1} \\
    residual complex & \bookref{sec:6.1}{VI.1} \\
    residually stable morphism & \bookref{sec:6.5}{VI.5} \\
    residue of a differential & \bookref{sec:7.1}{VII.1} \\
    residue symbol & \bookref{sec:3.9}{III.9} \\
    sheaf of families of supports & \bookref{sec:4.1}{IV.1} \\
    smooth morphism & \bookref{sec:3.1}{III.1} \\
    translation functor & \bookref{sec:1.1}{I.1} \\
    triangle & \bookref{sec:1.1}{I.1} \\
    triangulated category & \bookref{sec:1.1}{I.1} \\
    way-out functor & \bookref{sec:1.7}{I.7} \\
    $Z/Z^{\prime}$-skeleton & \bookref{sec:4.2}{IV.2} \\
\end{longtable}
\endgroup

\chapter{Index of Notations}
Of course there is the usual collection of variable notations: $A, B$ for abelian categories or rings, $F,G$ for functors or sheaves, $X,Y$ for preschemes, $x,y$ for points of them, $f,g$ for morphisms of preschemes, $\phi,\psi$ for morphisms of functors on sheaves, etc.


In this edition, ordinary $\Hom(F,G)$ denotes a group or module, while
$\sheafHom(F,G)$ and $\sheafTor(F,G)$ denote sheaf-valued functors.
Derived functors use upright $\RR$ and $\LL$: $\RHom(F^\bullet,G^\bullet)$
is global derived Hom, $\RsheafHom(F^\bullet,G^\bullet)$ is derived sheaf Hom,
and $F^\bullet\derivedtensor G^\bullet$ is the derived tensor product.
A bullet denotes a complex; in particular, $\Hom^\bullet$ and
$\sheafHom^\bullet$ denote Hom complexes when these are used explicitly.
The original edition used one underline for sheaves and two underlines
for derived objects; the notation here replaces that convention.

In the list of stable notations below, we have distinguished five categories: Latin alphabet, other alphabets, superscripts, subscripts, and arbitrary symbols.

\section*{Latin alphabet}
\begin{longtable}{@{}l p{8cm} l@{}}
	\toprule
	Symbol & Meaning & Reference \\
	\midrule
	\endfirsthead
	\midrule
	Symbol & Meaning & Reference \\
	\midrule
	\endhead
	\bottomrule
	\endfoot
	$\Ab$ & the category of abelian groups & \\
	acc & ascending chain condition & \hyperref[sec:5.2]{V.2} \\
	$\Ass$ & associated primes & \\
	$\B$ & boundaries of a complex & \\
	$\Coh$ & category of coherent sheaves & \hyperref[sec:2.1]{II.1} \\
	$\Coz$ & category of Cousin complexes & \hyperref[sec:4.3]{IV.3} \\
	$\mathscr{D}$ & sheaf-valued dualizing functor & \hyperref[sec:5.0]{V.0}, \hyperref[sec:7.3]{VII.3} \\
	$D$ & global dualizing functor & \hyperref[sec:7.3]{VII.3} \\
	$\DD(\mathcal{A})$ & derived category & \hyperref[sec:1.4]{I.4} \\
	$\DD(X)$ & $=\DD(\operatorname{Mod}(X))$ & \hyperref[sec:2.1]{II.1} \\
	dcc & descending chain condition & \hyperref[sec:5.2]{V.2} \\
	$\dim$ & dimension & \\
	$\operatorname{Dual}$ & category of dualizing complexes & \hyperref[sec:6.1]{VI.1} \\
	$\operatorname{E}$ & associated Cousin complex & \hyperref[sec:4.2]{IV.2}, \hyperref[sec:4.3]{IV.3} \\
	$\Ext$ & Ext groups & \hyperref[sec:1.6]{I.6} \\
	$\sheafExt$ & sheaf Ext & \hyperref[sec:2.3]{II.3} \\
	fid & finite injective dimension & \hyperref[prop:1.7.6]{I.7.6}, \hyperref[prop:2.7.20]{II.7.20} \\
	FR1--FR5 & axioms of multiplicative systems & \hyperref[sec:1.3]{I.3} \\
	$H$ & cohomology groups or modules & \hyperref[sec:1.1]{I.1}, \hyperref[sec:4.1]{IV.1} \\
	$\mathcal H$ & cohomology sheaves & \hyperref[sec:2.1]{II.1} \\
	$\Hom$ & Hom groups & \hyperref[sec:1.6]{I.6} \\
	$\sheafHom$ & sheaf Hom & \hyperref[sec:2.3]{II.3} \\
	$\RHom$ & global derived Hom & \hyperref[sec:1.6]{I.6}, \hyperref[sec:2.3]{II.3} \\
	$\RsheafHom$ & derived sheaf Hom & \hyperref[sec:2.3]{II.3} \\
	$i_x$ & sheaf functor associated to the local scheme at $x$ & \hyperref[sec:2.7]{II.7} \\
	$\operatorname{Icz}$ & category of injective Cousin complexes & \hyperref[sec:4.3]{IV.3} \\
	$\id$ & identity map & \\
	$\im$ & image & \\
	$J(x),J(x,x')$ & standard injectives & \hyperref[sec:2.7]{II.7} \\
	$\KK(\mathcal A)$ & homotopy category of complexes & \hyperref[sec:1.2]{I.2} \\
	$k(x)$ & residue field of a point $x$ & \\
	$K_\bullet(\underline{f}),K^\bullet(\underline{f};F)$ & Koszul complexes & \hyperref[sec:3.7]{III.7} \\
	$\ker$ & kernel & \\
	L1--L3 & axioms of direct systems & \hyperref[sec:1.3]{I.3} \\
	$\LL$ & left derived functor & \hyperref[sec:1.5]{I.5} \\
	lffr & locally free of finite rank & \hyperref[sec:2.5]{II.5} \\
	$\operatorname{Lno}$ & category of locally noetherian preschemes & \hyperref[sec:3.8]{III.8} \\
	$\operatorname{Mod}(X)$ & category of $\sO_X$-modules & \hyperref[sec:2.1]{II.1} \\
	$\Ob$ & objects of a category & \\
	pfid & pointwise finite injective dimension & \hyperref[sec:5.8]{V.8} \\
	$\operatorname{Ptwdual}$ & category of pointwise dualizing complexes & \hyperref[sec:6.1]{VI.1} \\
	$\operatorname{Q}$ & functor into derived category & \hyperref[sec:1.3]{I.3} \\
	$\Qco$ & category of quasi-coherent sheaves & \hyperref[sec:2.1]{II.1} \\
	$\operatorname{Qis}$ & quasi-isomorphisms & \hyperref[sec:1.4]{I.4} \\
	$\RR$ & right derived functor & \hyperref[sec:1.5]{I.5} \\
	$\operatorname{Res}$ & category of residual complexes & \hyperref[sec:6.1]{VI.1} \\
	$\operatorname{Res}$ & residue symbol & \hyperref[sec:3.9]{III.9} \\
	$\operatorname{Res}_x$ & residue of a differential & \hyperref[sec:7.1]{VII.1} \\
	$\Spec$ & spectrum of a ring & \\
	$\Supp$ & support of a sheaf or module & \hyperref[sec:2.7]{II.7} \\
	$\operatorname{T}$ & translation functor & \hyperref[sec:1.1]{I.1} \\
	$\sheafTor$ & sheaf Tor & \hyperref[sec:2.4]{II.4}  \\
	$\Tr$ & trace map & \hyperref[sec:3.10]{III.10}, \hyperref[sec:6.4]{VI.4}, \hyperref[sec:7.1]{VII.1} \\
	TR1--TR4 & axioms of a triangulated category & \hyperref[sec:1.1]{I.1} \\
	TRA & axioms of trace & \hyperref[thm:3.10.5]{III.10.5}, \hyperref[thm:6.4.2]{VI.4.2} \\
	tr.\ d & transcendence degree & \\
	$\operatorname{Trf}$ & trace for a finite morphism & \hyperref[sec:3.6]{III.6} \\
	$\operatorname{Trp}$ & trace for a projective morphism & \hyperref[sec:3.4]{III.4} \\
	VAR & axioms for $f^!$ or $f^\Delta$ & \hyperref[thm:3.8.7]{III.8.7}, \hyperref[thm:6.3.1]{VI.3.1} \\
	$\operatorname{Z}$ & cycles of a complex & \\
\end{longtable}

\section*{Greek and other alphabets}
\begin{longtable}{@{}l p{8cm} l@{}}
	\toprule
	Symbol & Meaning & Reference \\
	\midrule
	$\gamma$ & trace for smooth morphism & \hyperref[thm:3.3.4]{III.3.4}, \hyperref[cor:3.11.2]{III.11.2}, \hyperref[thm:7.4.1]{VII.4.1} \\
	$\Gamma$ & global sections of a sheaf & \\
	$\zeta$ & change of differentials & \hyperref[def:3.1.5]{III.1.5} \\
	$\eta$ & fundamental local isomorphism  & \hyperref[cor:3.7.3]{III.7.3} \\
	$\Theta,\underline{\Theta}$ & duality morphisms & \shortstack[l]{\hyperref[sec:3.5]{III.5}, \hyperref[sec:3.6]{III.6}, \hyperref[sec:3.11]{III.11},\\\hyperref[sec:5.6]{V.6}, \hyperref[sec:7.3]{VII.3}} \\
	$\xi$ & map of a functor to its derived functor & \hyperref[sec:1.5]{I.5} \\
	$\sigma_{>n}$ & truncated complex & \hyperref[sec:1.7]{I.7} \\
	$\tau_{>n}$ & truncated complex & \hyperref[sec:1.7]{I.7} \\
	$\psi$ & residue isomorphism & \hyperref[sec:3.8]{III.8} \\
	$\omega$ & sheaf of differentials & \hyperref[sec:3.1]{III.1} \\
	$\Omega$ & sheaf of differentials & \hyperref[sec:3.1]{III.1} \\
	$\sO_X$ & structure sheaf of a prescheme & \\
	$\mathfrak{m}$ & maximal ideal of a ring & \\
	$\mathfrak{p}$ & prime ideal of a ring & \\
	$\mathbb{A}$ & affine space & \\
	$\mathbb{P}$ & projective space & \\
	$\mathbb{Q}$ & rational numbers & \\
	$\mathbb{Z}$ & integers & \\
\end{longtable}

\section*{Subscripts}
\begin{longtable}{@{}l p{8cm} l@{}}
	\toprule
	Symbol & Meaning & Reference \\
	\midrule
	$\mathcal A'$ & as in $\KK_{\mathcal A'}(\mathcal{A}), \DD_{\mathcal A'}(\mathcal{A})$ & \hyperref[sec:1.4]{I.4} \\
	$c$ & as in $\DD_c(X)$ & \hyperref[sec:2.1]{II.1} \\
	fid & as in $\DD(X)_{\text{fid}}$ & \hyperref[sec:1.7]{I.7} \\
	fTd & as in $\DD(X)_{\text{fTd}}$ & \hyperref[sec:2.4]{II.4} \\
	$\operatorname{Gor}(Z^\bullet)$ & as in $\DD(X)_{\operatorname{Gor}(Z^\bullet)}$ & \hyperref[sec:4.3]{IV.3} \\
	$\mathfrak{p}$ & as in $A_\mathfrak{p},\,M_\mathfrak{p}$: localization at a prime ideal & \\
	qc & as in $\DD_{\text{qc}}(X)$ & \hyperref[sec:2.1]{II.1} \\
	$x$ & as in $\sO_x$: local ring at a point  $x$ & \\
	 & as in $\sF_x$: stalk of a sheaf at  $x$ & \\
	 & as in $\HH^i_x$: local cohomology & \\
	$Z$ & as in $\Gamma_Z,\HH^i_Z$ & \hyperref[sec:4.1]{IV.1} \\
	$\operatorname{I},\,\operatorname{II}$ & as in $\RR_{\operatorname{I}},\RR_{\operatorname{II}}$ & \hyperref[sec:1.6]{I.6} \\
	$>n$ & as in $\sigma_{>n},\tau_{>n}$ & \hyperref[sec:1.7]{I.7} \\
	$*$ & as in $f_*$: direct image & \hyperref[sec:2.2]{II.2} \\
\end{longtable}

\section*{Superscripts}
\begin{longtable}{@{}l p{8cm} l@{}}
	\toprule
	Symbol & Meaning & Reference \\
	\midrule
	$b$ & as in $\KK^b, \DD^b$: bounded & \hyperref[sec:1.2]{I.2}, \hyperref[sec:1.4]{I.4} \\
	$i$ & as in $X^i$:$i$-th term of a complex & \\
	& as in $H^i$:$i$-th cohomology group & \\
	& as in $\RR^i F$:$i$-th derived functor & \\
	$y$ & as in $f^y$ & \hyperref[sec:6.2]{VI.2} \\
	$z$ & as in $f^z$ & \hyperref[sec:6.2]{VI.2} \\
	$\bullet$ & as in $\sF^\bullet, X^\bullet$: denotes a complex  & \\
	$\bullet\bullet$ & as in $\mathcal{C}^{\bullet\bullet}$: double complex & \hyperref[sec:1.7]{I.7} \\
	$\mathrm{op}$ & as in $\KK(\mathcal A)^{\mathrm{op}}$: opposite category (a superscript bullet in the original) &\\
	$*$ & as in $f^*$: inverse image & \hyperref[sec:2.4]{II.4} \\
	$+$ & as in $\KK^+, \DD^+, \RR^+$: bounded below  & \hyperref[sec:1.2]{I.2}, \hyperref[sec:1.4]{I.4}, \hyperref[sec:1.5]{I.5} \\
	$-$ & as in $\KK^-, \DD^-, \RR^-$: bounded above  & \hyperref[sec:1.2]{I.2}, \hyperref[sec:1.4]{I.4}, \hyperref[sec:1.5]{I.5} \\
	$!$ & as in $f^!$ & \hyperref[chap:introduction]{Intr.}, \hyperref[sec:3.8]{III.8}, \hyperref[sec:7.3]{VII.3} \\
	$\flat$ & $f^\flat$ & \hyperref[sec:3.6]{III.6} \\
	$\sharp$ & $f^\sharp$ & \hyperref[sec:3.2]{III.2} \\
	$\Delta$ & $f^\Delta$ & \hyperref[sec:6.3]{VI.3} \\
	$\vee$ & as in $L^\vee$: dual sheaf & \hyperref[prop:2.5.16]{II.5.16}, \hyperref[sec:3.1]{III.1} \\
	$-$ & as in $\{x\}^-$: closure & \\
	$\sim$ & as in $\tilde{M}$: sheaf associated to a module & \\
	$\wedge$ & as in $\hat{A}$: completion of local ring & \\
\end{longtable}

\section*{Arbitrary Symbols}
\begingroup
\renewcommand{\arraystretch}{0.8}
\begin{longtable}{@{}l p{8cm} l@{}}
	\toprule
	Symbol & Meaning & Reference \\
	\midrule
	$\vert$ & restriction & \\
	$\gg$ & enough larger than & \\
	$\times$ & product & \\
	$\otimes$ & tensor product & \\
	$\derivedtensor$ & derived tensor product & \hyperref[sec:2.4]{II.4} \\
	$\oplus$ & direct sum & \\
	$\prod$ & product & \\
	$\coprod$ & disjoint union / sum & \\
	$[\ ]$ & shift operator & \hyperref[sec:1.2]{I.2} \\
	 & closed interval & \hyperref[sec:5.8]{V.8} \\
	& reference to Bibliography or other chapter & \hyperref[chap:introduction]{Intr.} \\
	$\bigcup$ & union & \\
	$\{\ \}$ & the set of & \\
	$\to$ & morphism & \\
	$\mapsto$ & effect of a map on elements & \\
	$\cdot\to$ & distinguished side of triangle & \hyperref[sec:1.1]{I.1} \\
	$\dashrightarrow$ & morphism being constructed & \\
	$\circ$ & composition of functors ($F\circ G$) & \hyperref[sec:1.2]{I.2} \\
\end{longtable}
\endgroup
\nocite{*}
\printbibliography[title={Bibliography}]
\appendix
\chapter{Cohomologie à Support Propre, et Construction du Foncteur $f^!$}
\markboth{Cohomologie \`a Support Propre}{Cohomologie \`a Support Propre}
par P.\ Deligne\footnote{Ceci est une version compl\'et\'ee d'une lettre de P.\ Deligne \`a R.\ Hartshorne (lettre du 3 Mars 1966). Les notes de bas de page ont \'et\'e ajout\'ees par le copiste.}

Verdier a montr\'e que dans le cas topologique, le formalisme de la dualit\'e de Poincar\'e se ramenait \`a des probl\`emes locaux en haut (voir \hyperref[app:bib1]{[1]}). Pour transposer sa construction au cadre sch\'ematique, il faut disposer d'une th\'eorie de la cohomologie ``\`a support propre'' pour les faisceaux coh\'erents.

Sauf mention explicite du contraire, tous les pr\'esch\'emas consid\'er\'es sont noeth\'eriens et les pr\'efaisceaux quasi-coh\'erents.

\section*{n\textdegree{} 1. Le sorite des pro-objets.}\label{app:sec1}

\begin{bookproposition}{1}\label{app:prop1}
Soit $\mathcal{C}$ une $\mathbb{U}$-cat\'egorie\footnote{$\mathbb{U}$ d\'esigne un univers fix\'e dans toute la suite.} o\`u existent les limites inductives finies. Soit $h$ un foncteur $\mathcal{C}^{\circ} \to (\text{ens})$. Les conditions suivantes sont \'equivalentes:
\begin{enumerate}
\item[(i)] $h$ est limite inductive, selon un petit\footnote{``petit'' = ``$\in \mathbb{U}$''.} ensemble ordonn\'e filtrant, de foncteurs repr\'esentables.
\item[(ii)] $h$ est limite inductive, selon une petite cat\'egorie filtrante, de foncteurs repr\'esentables.
\item[(iii)] $h$ transforme $\varinjlim$ finies en $\varprojlim$ finies, et il existe un petit ensemble d'objets tel que tout \'el\'ement d'un $h(X)$ se factorise par l'un d'eux.
\end{enumerate}
\end{bookproposition}

Les implications (i) $\Rightarrow$ (ii) $\Rightarrow$ (iii) sont triviales et (iii) $\Rightarrow$ (ii) standard ($h$ est limite des $h_F$ selon la cat\'egorie des $(F, \alpha)$ pour $\alpha \in h(F)$ et $F$ dans la sous-cat\'egorie de $\mathcal{C}$ stable par $\varinjlim$ finie engendr\'ee par le petit ensemble donn\'e). Reste \`a prouver que (ii) $\Rightarrow$ (i). Si $\mathcal{L}$ et $\mathcal{M}$ sont deux cat\'egories filtrantes, un foncteur $F: \mathcal{L} \to \mathcal{M}$ est dit cofinal si $\varinjlim_{L\in\mathcal L} h_{F(L)}$ est le foncteur final. Pour tout $G\colon\mathcal M\to\mathcal N$, on a alors $\varinjlim_{\mathcal M}G=\varinjlim_{\mathcal L}(G\circ F)$.

\par\medskip\noindent\textbf{Lemme.} Pour toute petite cat\'egorie filtrante $\mathcal{L}$, il existe un foncteur cofinal d'un petit ensemble ordonn\'e filtrant dans $\mathcal{L}$.

La premi\`ere projection: $\mathcal{L} \times \mathbb{N} \longrightarrow \mathcal{L}$, est un foncteur cofinal ($\mathbb{N}$ muni de l'ordre naturel); ceci permet de se ramener \`a supposer que $\operatorname{Ob}\mathcal{L}$, pr\'eordonn\'e par $\Hom(X,Y) \neq \emptyset$, n'a pas de plus grand \'el\'ement. On ordonne par inclusion l'ensemble $E$ des sous-cat\'egories finies de $\mathcal{L}$ ayant un seul objet final (fini signifie d'ensemble de fl\`eches fini). On d\'efinit un foncteur de $E$ dans $\mathcal{L}$ en associant \`a chacune de ces cat\'egories son objet final. Sous les hypoth\`eses faites, $E$ est filtrant et ce foncteur cofinal.

Les foncteurs v\'erifiant les conditions \'equivalentes de la proposition 1 sont les \emph{Ind-objets} de $\mathcal{C}$. Ils forment une sous-cat\'egorie pleine $\operatorname{Ind}\mathcal{C}$ de $\operatorname{Fun}(\mathcal C^{\mathrm{op}},\mathrm{Ens})$, stable par limite inductive. $\mathcal{C} \longmapsto \operatorname{Ind}\mathcal{C}$ est un foncteur $(\text{cat}) \to (\text{cat})$. Si les limites inductives filtrantes existent dans $\mathcal{C}$, pour tout $h \in \operatorname{Ind}\mathcal{C}$, le foncteur $X\longmapsto\Hom(h,h_X)$ est corepr\'esentable, d'o\`u un foncteur, not\'e $\varinjlim$, de $\operatorname{Ind}\mathcal{C}$ dans $\mathcal{C}$. En particulier, on a toujours un foncteur $\operatorname{Ind}\operatorname{Ind}\mathcal{C} \longrightarrow \operatorname{Ind}\mathcal{C}$, et tout foncteur $\mathcal C\longrightarrow\operatorname{Ind}\mathcal D$ se prolonge en un foncteur $\operatorname{Ind}\mathcal C\longrightarrow\operatorname{Ind}\mathcal D$. Dans $\operatorname{Ind}\mathcal{C}$, les limites inductives filtrantes sont exactes, et seront not\'ees ``$\varinjlim$''; en particulier, identifiant $\mathcal{C}$ \`a une sous-cat\'egorie pleine de $\operatorname{Ind}\mathcal{C}$, si $(X_i)_{i \in I}$ est un syst\`eme inductif filtrant dans $\mathcal{C}$, $\varinjlim_{i\in I}X_i=\varinjlim\bigl(\text{``}\varinjlim_{i\in I}\text{''}X_i\bigr)$.

Pour qu'un foncteur sur $\operatorname{Ind}\mathcal{C}$ soit repr\'esentable, il faut et suffit qu'il transforme $\varinjlim$ finies en $\varprojlim$ finies, $\varinjlim$ filtrantes en $\varprojlim$ filtrantes et que sa restriction \`a $\mathcal{C}$ satisfasse \`a la condition de petitesse de la prop.\ 1 (iii) (en effet, la proposition 1 montre alors que sa restriction \`a $\mathcal{C}$ est un Ind-objet, auquel il est partout \'egal vu la condition sur les limites).

Ce qui pr\'ec\`ede, sauf la prop.\ 1 (iii) et l'assertion pr\'ec\'edente reste vrai en utilisant partout des limites de suites (ou d\'enombrables, c'est la m\^eme chose).

On d\'efinit par dualit\'e la cat\'egorie $\operatorname{pro}\mathcal{C}$ des pro-objets (sous-cat\'egorie pleine de $\operatorname{Fun}(\mathcal C,\mathrm{Ens})^{\mathrm{op}}$).

\begin{bookproposition}{2}\label{app:prop2}
Soit $X$ un pr\'esch\'ema quasi-compact quasi-s\'epar\'e (non n\'ecessairement noeth\'erien). La cat\'egorie des faisceaux quasi-coh\'erents sur $X$ est \'equivalente \`a la cat\'egorie des Ind-objets de la cat\'egorie des faisceaux quasi-coh\'erents de pr\'esentation finie sur $X$.
\end{bookproposition}

La fl\`eche est $\text{``}\varinjlim\text{''}\mathcal F_i\longmapsto\varinjlim_i\mathcal F_i$. Si $\mathcal F$ est de p.f.\ (pr\'esentation finie), pour tout syst\`eme inductif filtrant $(\mathcal G_i)_{i\in I}$, on a
\[
\Hom_X\bigl(\mathcal F,\varinjlim_i\mathcal G_i\bigr)
\cong\varinjlim_i\Hom_X(\mathcal F,\mathcal G_i)
\] (par un recollement fini qui commute \`a $\varinjlim$, on se ram\`ene au cas affine). Le foncteur pr\'ec\'edent est donc pleinement fid\`ele. Son image est stable par limites inductives filtrantes et sommes finies. Il suffit de prouver que pour tout $x\in X$, tout faisceau quasi-coh\'erent $\mathcal F$ sur $X$ et tout $s\in\mathcal F_x$, il existe un faisceau $\mathcal G$ de p.f.\ et une fl\`eche $\mathcal G\to\mathcal F$ dont l'image contient $s$: $\mathcal F$ sera limite d'images de faisceaux de p.f. et chacune d'elles quotient d'un faisceau de p.f.\ par une limite de sous-faisceaux de type fini.

Soit sur un voisinage quasi-compact $U$ de $x$ une fl\`eche $f: \mathcal{C} \longrightarrow \mathcal{F}$, $\mathcal{C}$ de p.f., $s$ dans l'image. Il suffit de prolonger $\mathcal{C}$ et $f$ \`a $X$ entier, et, proc\'edant pas \`a pas, il suffit de savoir prolonger \`a $U \cup V$ si $V$ est affine, ce qui revient \`a effectuer un prolongement de $U \cap V$ \`a $V$:

\par\medskip\phantomsection\label{app:extension-lemma}\noindent\textbf{Lemme.}
Soit $\mathcal{C}$ un faisceau de p.f.\ sur un ouvert quasi-compact $U$ d'un sch\'ema affine $X$, $\mathcal{F}$ un faisceau sur $X$ et $f$ une fl\`eche de $\mathcal{C}$ dans $\mathcal{F}|_U$. Il est possible de prolonger $\mathcal{C}$ et $f$ en $\overline{\mathcal{C}}$ et $\overline{f}$, d\'efinis sur $X$ ($\overline{\mathcal{C}}$ \'etant de p.f.).
\par\medskip

Soit $\iota$ l'inclusion de $U$ dans $X$ et $\mathcal{C}_1$ le produit fibr\'e de $\iota_*\mathcal{C}$ et $\mathcal{F}$ sur $\iota_*\iota^*\mathcal{F}$. $\mathcal{C}_1$ prolonge $\mathcal{C}$ et s'envoie dans $\mathcal{F}$. Si on repr\'esente $\mathcal{C}_1$ comme limite de ses sous-faisceaux de type fini, comme $U$ est quasi-compact et $\mathcal{C}$ de p.f., l'un d'eux prolonge d\'ej\`a $\mathcal{C}$. Si on repr\'esente ce dernier comme quotient de $\mathcal{O}^n$ par une limite de sous-faisceaux de type fini de $\mathcal{O}^n$, on voit de m\^eme qu'on peut remplacer $\mathcal{C}_1$ par un faisceau de p.f.

\begin{bookcorollary}{1}\label{app:cor1}
Pour qu'un foncteur contravariant additif $F: (\text{q coh}) \to \Ab$ soit repr\'esentable, il faut et suffit qu'il soit exact \`a gauche et transforme $\varinjlim$ filtrante en $\varprojlim$ filtrante.
\end{bookcorollary}

\begin{bookcorollary}{2}\label{app:cor2}
Tout faisceau de p.f.\ d\'efini sur un ouvert quasi-compact $U$ de $X$ se prolonge en un faisceau de p.f.\ sur $X$.
\end{bookcorollary}

Soit $\mathcal{A}$ une cat\'egorie ab\'elienne. On aura \`a travailler de fa\c{c}on essentielle dans $\operatorname{pro}\DD^b\mathcal{A}$, sous-cat\'egorie pleine de $\operatorname{pro}\DD(\mathcal{A})$ form\'ee des ``$\varprojlim$'' $K_i$ o\`u la cohomologie de $K_i$ est uniform\'ement born\'ee avec $i$. $H^p$ se prolonge en un foncteur $\operatorname{pro}\DD(\mathcal{A}) \longrightarrow \operatorname{pro}\mathcal{A}$.

\begin{bookproposition}{3}\label{app:prop3}
Les foncteurs $H^p$ forment un syst\`eme conservatif\footnote{i.e.\ si $u$ est une fl\`eche telle que les $H^p(u)$ soient inversibles, $u$ est inversible.} de foncteurs sur $\operatorname{pro}\DD^b\mathcal{A}$.
\end{bookproposition}

Soit $f: K \to L$ une fl\`eche de $\operatorname{pro}\DD^b\mathcal{A}$ telle que les $H^p(K) \to H^p(L)$ soient des isomorphismes (dans $\operatorname{pro}\mathcal{A}$). Il s'agit de v\'erifier que pour tout $M$ dans $\DD^b(\mathcal{A})$, $\Hom(K,M) \longleftarrow \Hom(L,M)$ est un isomorphisme. Posons $K = \text{``$\varprojlim$''} K_{\kappa}$. On dispose d'une suite spectrale convergente\footnote{Correction \'editoriale d'indices: le texte original porte $H^qK_{\kappa}$. Pour l'aboutissement en degr\'e $p+q$, il faut $H^{-q}K_{\kappa}$.}
\[
\Ext^p(H^{-q}K_{\kappa},M)\Longrightarrow\Ext^{p+q}(K_{\kappa},M),
\]
d'o\`u, par passage \`a la limite,
\[
\varinjlim_{\kappa}\Ext^p(H^{-q}K_{\kappa},M)
\Longrightarrow\varinjlim_{\kappa}\Ext^{p+q}(K_{\kappa},M),
\]
encore convergente car $K$ est uniform\'ement born\'e. Utilisant que tout
foncteur (ici $\Ext$) passe aux pro-objets et le foncteur
$\varinjlim\colon\operatorname{Ind}\Ab\to\Ab$, on peut r\'ecrire
\[
\varinjlim\Ext^p(H^{-q}K,M)
\Longrightarrow\varinjlim\Ext^{p+q}(K,M)
=\Hom(K[-p-q],M).
\]
Cette suite spectrale reste d'ailleurs vraie pour $M \in \DD(\mathcal{A})$. La proposition r\'esulte aussit\^ot de l'existence de ces suites.

\section*{n\textdegree{} 2. Lemmes fondamentaux.}

Rappelons que les pr\'esch\'emas consid\'er\'es sont noeth\'eriens.

\begin{bookproposition}[Th\'eor\`eme de dualit\'e pour une immersion ouverte]{4}\label{app:prop4}
Soient $U$ un ouvert de $X$, $j$ la fl\`eche d'inclusion, $\mathcal{J}$ un id\'eal d\'efinissant le ferm\'e compl\'ementaire, $\mathcal{F}$ un faisceau coh\'erent sur $U$ et $\mathcal{G}$ quasi-coh\'erent sur $X$. Soit $\overline{\mathcal{F}}$ un prolongement coh\'erent de $\mathcal{F}$. On a\footnote{Cette formule est \'evidemment bien connue, cf.\ p.\ ex.\ la th\`ese de P.\ Gabriel.}
\[
\Hom_{U}(\mathcal{F}, j^{*}\mathcal{G}) = \Hom_{X}(\text{``}\varprojlim_n\text{''} \mathcal{J}^{n}\overline{\mathcal{F}}, \mathcal{G}).
\]
\end{bookproposition}

La fl\`eche $\varinjlim \Hom_X(\mathcal{J}^{n}\overline{\mathcal{F}},\mathcal{G}) \longrightarrow \Hom_U(\mathcal{F},\mathcal{G})$ est injective: si l'image par $f$ de $\mathcal{J}^{n}\overline{\mathcal{F}}$ dans $\mathcal{G}$ a son support dans le compl\'ementaire de $U$, elle est annul\'ee par une puissance de $\mathcal{J}$, soit $\mathcal{J}^{k}$, et l'image de $\mathcal{J}^{n+k}\overline{\mathcal{F}}$ est nulle.

Supposons maintenant $X$ affine et soit $f \in \Hom_U(\mathcal{F},\mathcal{G})$. Pour $k$ assez grand, toute section $s$ de $\mathcal{J}^{k}\overline{\mathcal{F}}$ sur $X$ a une image (sur $U$) qui se prolonge en une section de $\mathcal{G}$ sur $X$. Rempla\c{c}ant $\overline{\mathcal{F}}$ par $\mathcal{J}^{k}\overline{\mathcal{F}}$, on peut supposer disposer d'un diagramme
\[
\begin{tikzcd}
0 \ar[r] & R \ar[r] \ar[d,dashed] & \mathcal{O}^{n} \ar[r] \ar[d,equals] & \overline{\mathcal F} \ar[r] \ar[d,dashed] & 0 \\
0 \ar[r] & R_1 \ar[r] & \mathcal{O}^{n} \ar[r] & \mathcal G
\end{tikzcd}
\]
o\`u les fl\`eches pointill\'ees sont d\'efinies sur $U$. Pour $\ell$ assez grand, en vertu de Krull et du fait que $R/R \cap R_1$ est concentr\'e hors de $U$, on a $R \cap \mathcal{J}^{\ell} \cdot \mathcal{O}^{n} \subset R_1$, ce qui permet de prolonger $f$ en $\overline{f}: \mathcal{J}^{\ell}\overline{\mathcal{F}} \to \mathcal{G}$.

Dans le cas g\'en\'eral, soient $U_i$ un recouvrement affine fini de $X$ et $U_{ijk}$ des recouvrements affins finis des $U_{ij} = U_i \cap U_j$. Les $\varinjlim$ filtrantes commutent aux produits finis. Posons, les restrictions aux ouverts concern\'es \'etant sous-entendues,
\[
\begin{aligned}
L(V)&=\Hom_{U\cap V}(\mathcal F,j^*\mathcal G),\\
P(V)&=\varinjlim_n\Hom_V(\mathcal J^n\overline{\mathcal F},\mathcal G).
\end{aligned}
\]
On obtient le diagramme \`a lignes exactes
\[
\begin{tikzcd}[column sep=large,row sep=large]
0 \ar[r] & L(X) \ar[r]
 & \prod_i L(U_i) \ar[r,shift left=.5ex] \ar[r,shift right=.5ex]
 & \prod_{i,j,k}L(U_{ijk}) \\
0 \ar[r] & P(X) \ar[r] \ar[u]
 & \prod_i P(U_i) \ar[r,shift left=.5ex] \ar[r,shift right=.5ex] \ar[u,"\sim"']
 & \prod_{i,j,k}P(U_{ijk}) \ar[u,"\sim"']
\end{tikzcd}
\]
ce qui ach\`eve la d\'emonstration.

La proposition~\ref{app:prop4} donne une formule explicite pour le foncteur
$(\text{coh sur }U)\to\operatorname{pro}(\text{coh sur }X)$
``adjoint'' \`a gauche de $j^*$. On notera ce foncteur $j_!$
(``prolongement par z\'ero'')\footnote{Cette terminologie et notation conflictent \'evidemment avec ceux g\'en\'eralement admis (livre de Godement, SGA 1962, \dots).};
il r\'esulte de Krull qu'il est exact; cela r\'esulte aussi de ce que,
$X$ \'etant noeth\'erien, un injectif de la cat\'egorie des faisceaux
quasi-coh\'erents sur $X$, restreint \`a $U$, reste injectif quasi-coh\'erent sur $U$.

\begin{bookproposition}[Ind\'ependance de la compactification]{5}\label{app:prop5}
Soit $f\colon X\to Y$ un morphisme propre, induisant un isomorphisme
entre l'ouvert $U$ de $Y$ et $f^{-1}(U)$:
\[
\begin{tikzcd}
U \ar[r,hook] \ar[d,equals] & X \ar[d,"f"] \\
U \ar[r,hook] & Y.
\end{tikzcd}
\]
Soit $\mathcal J$ un faisceau d'id\'eaux d\'efinissant $Y\setminus U$,
et soit $\mathcal J'=\mathcal J\mathcal O_X$ l'id\'eal induit sur $X$.
Si $\mathcal F$ est coh\'erent sur $X$, on a
\begin{enumerate}
\item[(i)] si $k>0$, le syst\`eme projectif
$\RR^k f_*(\mathcal J'{}^n\mathcal F)$ est essentiellement nul
(i.e.\ d\'efinit le pro-objet nul);
\item[(ii)] si $k=0$, pour $n$ assez grand,
\[
f_*(\mathcal J'{}^{n+1}\mathcal F)
=\mathcal J\cdot f_*(\mathcal J'{}^n\mathcal F).
\]
\end{enumerate}
\end{bookproposition}

On a
\[
\bigoplus_{n\ge0}\RR^k f_*(\mathcal J'{}^n\mathcal F)
\quad\text{un }\bigoplus_{n\ge0}\mathcal J^n\text{-module de type fini}
\quad\text{\cite[\nopp III]{EGA}}.
\]
Donc, pour $n$ assez grand,
\[
\mathcal J\otimes\RR^k f_*(\mathcal J'{}^n\mathcal F)
\longrightarrow\RR^k f_*(\mathcal J'{}^{n+1}\mathcal F)
\]
est surjectif; (ii) en r\'esulte. Si $k>0$, $\RR^k f_*$ est nul sur $U$;
il existe $N$ tel que $\mathcal J^N\RR^k f_*(\mathcal J'{}^n\mathcal F)$
soit nul (quel que soit $n$).
Pour all\'eger le diagramme, posons $M_n=\RR^k f_*(\mathcal J'{}^n\mathcal F)$.
Le diagramme
\[
\begin{tikzcd}[row sep=large,column sep=large]
\mathcal J^p\otimes M_n \ar[r] \ar[d]
 & M_{n+p} \ar[r] \ar[d] & 0\quad(n\ge n_0) \\
\mathcal J^p\cdot M_n \ar[r] & M_n
\end{tikzcd}
\]
prouve que pour $N$ assez grand et $k>0$,
\[
\RR^k f_*(\mathcal J'{}^{n+N}\mathcal F)
\longrightarrow\RR^k f_*(\mathcal J'{}^n\mathcal F)
\]
est nul, soit (i).

\section*{n\textdegree{} 3. D\'efinition de $\RR f_{!}$.}\label{app:sec3}

Un morphisme $f$ (resp.\ un couple de morphismes composables, resp.\ un triple \dots) sera dit \emph{compactifiable} si on peut respectivement trouver des diagrammes commutatifs
\[
\begin{tikzcd}[row sep=large,column sep=large]
X \ar[r,hook] \ar[d,"f"'] & \overline X \ar[dl] \\
S
\end{tikzcd}
\qquad
\begin{tikzcd}[row sep=large,column sep=large]
X \ar[r,hook] \ar[d,"g"'] & \overline X \ar[r,hook] \ar[dl]
 & \overline{\overline X} \ar[dl] \\
Y \ar[r,hook] \ar[d,"f"'] & \overline Y \ar[dl] \\
S
\end{tikzcd}
\]
\[
\begin{tikzcd}[row sep=large,column sep=large]
X \ar[r,hook] \ar[d,"h"'] & \overline X \ar[r,hook] \ar[dl]
 & \overline{\overline X} \ar[r,hook] \ar[dl]
 & \overline{\overline{\overline X}} \ar[dl] \\
Y \ar[r,hook] \ar[d,"g"'] & \overline Y \ar[r,hook] \ar[dl]
 & \overline{\overline Y} \ar[dl] \\
Z \ar[r,hook] \ar[d,"f"'] & \overline Z \ar[dl] \\
S
\end{tikzcd}
\]
o\`u les fl\`eches horizontales sont des immersions ouvertes et les fl\`eches obliques sont propres. Tout morphisme compactifiable est s\'epar\'e de type fini, et Nagata affirme dans \hyperref[app:bib2]{[2]} que sa d\'emonstration prouve la r\'eciproque pour les sch\'emas noeth\'eriens int\`egres, mais les hypoth\`eses qu'il fait sur la base ne sont pas claires.\footnote{Mumford aurait v\'erifi\'e via la d\'emonstration de Nagata que tout morphisme s\'epar\'e de type fini des sch\'emas noeth\'eriens est compactifiable.}

On se propose, pour tout $f$ compactifiable, de d\'efinir $\RR f_{!}: \operatorname{pro}\DD_{\operatorname{coh}}^{b}(X) \longrightarrow \operatorname{pro}\DD_{\operatorname{coh}}^{b}(S)$. Rappelons que la cat\'egorie $\DD_{\operatorname{coh}}^{b}(S)$, cat\'egorie d\'eriv\'ee born\'ee de la cat\'egorie des faisceaux coh\'erents sur $S$, est sous-cat\'egorie pleine de la cat\'egorie d\'eriv\'ee de celle de tous les faisceaux (non n\'ecessairement quasi-coh\'erents) sur $S$, l'image essentielle \'etant form\'ee des complexes \`a cohomologie coh\'erente et born\'ee.

Si $f$ est propre, on dispose de $\RR f_{*}\colon \DD_{\operatorname{coh}}^{b}(X) \longrightarrow \DD_{\operatorname{coh}}^{b}(S)$, de dimension finie, qui induit $\RR f_{!}: \operatorname{pro}\DD_{\operatorname{coh}}^{b}(X) \longrightarrow \operatorname{pro}\DD_{\operatorname{coh}}^{b}(S)$. Si $f$ est une immersion ouverte, $f_{!}: (\text{coh sur } X) \longrightarrow \operatorname{pro}(\text{coh sur } S)$ induit $\DD^{b}(\text{coh sur } X) \longrightarrow \DD^{b}\operatorname{pro}(\text{coh sur } S)$ qui s'envoie dans $\operatorname{pro}\DD_{\operatorname{coh}}^{b}(S)$, et cette fl\`eche se prolonge en $\RR f_{!}: \operatorname{pro}\DD_{\operatorname{coh}}^{b}(X) \longrightarrow \operatorname{pro}\DD_{\operatorname{coh}}^{b}(S)$. Si on prolonge le complexe coh\'erent $K$ sur $X$ en $\overline{K}$ sur $S$, son image est $\text{``}\varprojlim_n\text{''}\mathcal J^n\overline K$, o\`u $\mathcal{J}$ d\'efinit $S \setminus X$.

Pour un compos\'e $f = gh$ ($g$ propre, $h$ immersion ouverte) on prend $\RR f_{!} = \RR g_{!}\RR h_{!}$. Il faut v\'erifier:

\subsection*{Ind\'ependance de la compactification.}

Consid\'erons un diagramme commutatif
\[
\begin{tikzcd}[row sep=large,column sep=large]
X \ar[r,hook,"j'"] \ar[dr,hook,"j''"'] \ar[dd,"f"']
 & X' \ar[d,"g"] \\
 & X'' \ar[dl,"\overline f"] \\
S
\end{tikzcd}
\]
o\`u $j'$ et $j''$ sont des immersions ouvertes, $\overline f$ et $g$
des applications propres. Il faut prouver que
\[
\RR\overline f_*\RR j''_!
=\RR(\overline f g)_*\RR j'_!.
\]
Puisque $\RR(\overline f g)_*=\RR\overline f_*\RR g_*$, il suffit de prouver
\[
\RR j''_!\cong\RR g_*\RR j'_!.
\]
On se ram\`ene \`a supposer $X$ dense dans $X'$, pour \^etre dans les
conditions de la Proposition~\ref{app:prop5}.
Soit $K$ un complexe coh\'erent born\'e sur $X$, prolong\'e en $\overline K$
sur $X'$. Le complexe $g_*\overline K$ prolonge $K$ sur $X''$, et on a une fl\`eche
\begin{equation}\tag{1}\label{app:eq:independence}
g_*(\mathcal J^n\overline K)\longrightarrow\RR g_*(\mathcal J^n\overline K).
\end{equation}
Ici $\mathcal J$ d\'efinit $X'\setminus X$. On peut le choisir induit
d'un id\'eal d\'efinissant $X''\setminus X$; c'est le choix utilis\'e pour
appliquer la Proposition~\ref{app:prop5}.
La partie (ii) de cette proposition donne, dans la cat\'egorie des pro-objets,
\[
j''_!K\cong\text{``}\varprojlim_n\text{''}\,g_*(\mathcal J^n\overline K).
\]
Par passage \`a la limite projective formelle, on trouve la suite spectrale
\[
\text{``}\varprojlim_n\text{''}\,\RR^p g_*(\mathcal J^n\overline K^q)
\Longrightarrow
\text{``}\varprojlim_n\text{''}\,\RR^{p+q}g_*(\mathcal J^n\overline K).
\]
La Proposition~\ref{app:prop5}(i) montre qu'elle d\'eg\'en\`ere: les termes
de degr\'e $p>0$ sont des pro-objets nuls. Il r\'esulte alors de la
Proposition~\ref{app:prop3} que la fl\`eche d\'eduite de
\eqref{app:eq:independence} par passage \`a la limite projective formelle
est un isomorphisme, d'o\`u la formule voulue.

\subsection*{Il reste \`a v\'erifier}

\begin{enumerate}
\item[(a)] pour tout diagramme, on a une compatibilit\'e. Notons $X'''$ la troisi\`eme compactification:
\[
\begin{tikzcd}[row sep=large,column sep=large]
X \ar[r,hook,"j'"] \ar[dr,hook,"j''"] \ar[ddr,hook,"j'''"'] \ar[ddd,"f"']
 & X' \ar[d,"h"] \\
 & X'' \ar[d,"g"] \\
 & X''' \ar[dl,"\overline f"] \\
S
\end{tikzcd}
\]
Les identifications successives doivent \^etre compatibles:
\[
\begin{aligned}
\RR\overline f_*\RR j'''_!
&\cong\RR\overline f_*\RR g_*\RR j''_!\\
&\cong\RR\overline f_*\RR g_*\RR h_*\RR j'_!\\
&\cong\RR\overline f_*\RR(gh)_*\RR j'_!\\
&\cong\RR\overline f_*\RR j'''_!.
\end{aligned}
\]
Comme deux compactifications peuvent toujours \^etre coiff\'ees par une
troisi\`eme (par exemple $X'\times_S X''$), ceci suffit \`a prouver
l'ind\'ependance de l'arbitraire.
\item[(b)] pour tout couple compactifiable, une identification
\[
\RR(fg)_!\cong\RR f_!\RR g_!.
\]
\item[(c)] pour tout triple compactifiable, une compatibilit\'e:
\[
\begin{tikzcd}[row sep=large,column sep=large]
\RR(fgh)_! \ar[r,"\sim"] \ar[d,"\sim"']
 & \RR(fg)_!\RR h_! \ar[d,"\sim"] \\
\RR f_!\RR(gh)_! \ar[r,"\sim"']
 & \RR f_!\RR g_!\RR h_!.
\end{tikzcd}
\]
\end{enumerate}

On aurait alors prouv\'e

\begin{booktheorem}{1}\label{app:thm1}
Pour $f: X \longrightarrow Y$ compactifiable, $K \in \operatorname{pro}\DD_{\operatorname{coh}}^{b}(X)$, $L \in \operatorname{pro}\DD_{\operatorname{coh}}^{b}(Y)$ posons $\Hom_{f}(K,L) = \Hom_{Y}(\RR f_{!}K,L)$. Modulo des questions de compactifiabilit\'e, on fait ainsi des cat\'egories $\operatorname{pro}\DD_{\operatorname{coh}}^{b}(X)$ une cat\'egorie cofibr\'ee sur les pr\'esch\'emas noeth\'eriens (les fl\`eches \'etant compactifiables \dots).
\end{booktheorem}

\section*{n\textdegree{} 4. D\'efinition de $\RR f^{!}$.}\label{app:sec4}

Il r\'esulte d'un tapis g\'en\'eral de Verdier \hyperref[app:bib1]{[1]} que le foncteur $\RR f_{*}: \DD(X) \longrightarrow \DD(S)$ a un adjoint \`a \emph{droite} $\DD^{+}(S) \longrightarrow \DD^{+}(X)$; il ne m\'erite un nom, soit $\RR f^{!}$, que pour $f$ propre, et sa d\'efinition ne se pr\^ete pas directement au calcul.

Il faut tout d'abord expliciter un proc\'ed\'e de calcul de $\RR f_{*}$ au niveau des complexes, du type $\RR f_{*}K = f_{*}C^{*}(K)$ o\`u $C^{*}$ est une r\'esolution acyclique finie d\'ependant de fa\c{c}on fonctorielle, exacte, et compatible avec les limites inductives filtrantes de l'objet auquel on l'applique. \`A ce moment, pour tout injectif $I$ sur $S$, $\Hom(f_{*}C^{p}\mathcal{F}, I)$ est exact en $\mathcal{F}$ quasi-coh\'erent sur $X$, et transforme $\varinjlim$ filtrantes en $\varprojlim$ filtrantes; il est donc repr\'esentable par $f_{p}^{!}I$, injectif quasi-coh\'erent. Les $f_{p}^{!}I$ forment un complexe, ce qui d\'efinit $\RR f^{!}:\DD^{+}(S)\longrightarrow\DD^{+}(X)$, adjoint \`a droite \`a $\RR f_{*}:\DD(X)\longrightarrow\DD(S)$.

Voici comment d\'efinir une telle r\'esolution:
\begin{enumerate}
\item[1)] Si $S$ est s\'epar\'e: on prend un recouvrement fini de $X$ par des ouverts affines sur $S$ (par exemple affines) et $\RR f_{*}K = f_{*}\check{C}(\mathfrak{U},K)$ (complexe de \v{C}ech altern\'e).

\item[2)] Sinon: les r\'esolutions flasques canoniques ont ici deux d\'efauts: (a) $f_{*}C^{*}$ n'est plus quasi-coh\'erent; cela ne porte pas \`a cons\'equence tant que, comme ici, les injectifs quasi-coh\'erents sont injectifs en tant que faisceaux. (b) $C^{*}$ ne commute pas aux limites inductives filtrantes. Il est facile de corriger ce d\'efaut: on repr\'esente un quasi-coh\'erent comme $\mathcal{F} = \varinjlim \mathcal{F}_i$ ($\mathcal{F}_i$ de pr\'esentation finie) et on pose $C'{}^{*}\mathcal{F} = \varinjlim C^{*}\mathcal{F}_i$. Cela ne d\'epend pas de l'arbitraire.
\end{enumerate}

Pour une immersion ouverte, on prendra $\RR f^{!}K = f^{*}K$. Pour un morphisme compos\'e $f = gh$ ($g$ propre, $h$ immersion ouverte), on prendra $\RR f^{!} = \RR h^{!}\RR g^{!}$\footnote{Correction \`a l'original: pour $X\xrightarrow{h}\overline{X}\xrightarrow{g}S$, les foncteurs contravariants se composent dans l'ordre $\RR h^{!}\RR g^{!}$.}. Mettant bout \`a bout la \bookref{app:prop4}{proposition 4} et ce qui pr\'ec\`ede, on trouve

\begin{booktheorem}{2}\label{app:thm2}
(Dualit\'e) Soit $f: X \longrightarrow S$ compactifiable, $f = gh$. Les foncteurs $\RR f_{!}: \operatorname{pro}\DD_{\operatorname{coh}}^{b}(X) \longrightarrow \operatorname{pro}\DD_{\operatorname{coh}}^{b}(S)$ et $\RR f^{!}: \DD^{+}(S) \longrightarrow \DD^{+}(X)$ sont ``adjoints'' l'un de l'autre.
\end{booktheorem}

Il faudrait v\'erifier l'ind\'ependance de la compactification et un sorite de compatibilit\'es et d'identifications analogue au cas de $\RR f_{!}$. Des cas particuliers du formulaire r\'esultent de la formule d'adjonction (unicit\'e d'un vrai adjoint \dots).

On pourra dans le cas g\'en\'eral s'appuyer sur la proposition suivante:

\begin{bookproposition}{6}\label{app:prop6}
Soit $f: X \longrightarrow S$ compactifiable et $L \in \DD_{\operatorname{qc}}^{+}(S)$. La fibre en $x \in X$ de $\mathcal{H}^{p}\RR f^{!}L$ est donn\'ee par
\[
(\mathcal{H}^{p}\RR f^{!}L)_{x} = \varinjlim_{x \in U} \Hom_{S}^{p}(\RR f_{!}(U,\mathcal{O}_{U}), L).
\]
\end{bookproposition}

On peut supposer $X$ propre et $L$ donn\'e comme complexe de faisceaux injectifs quasi-coh\'erents. Soient $U$ un ouvert affine de $X$ et $\mathcal{J}$ un id\'eal qui d\'efinit $X\setminus U$.
\begin{align*}
(\mathcal{H}^{p}\RR f^{!}L)(U)
 &= H^{p}\bigl((\RR f^{!}L)(U)\bigr) \\
 &= H^{p}\Hom_{\mathcal{O}_{U}}^{\bullet}
       (\mathcal{O}_{U},(\RR f^{!}L)|_{U}) \\
 &= H^{p}\varinjlim_{n}\Hom_{X}^{\bullet}
       (\mathcal{J}^{n},\RR f^{!}L) \\
 &= \varinjlim_{n}\Hom_{S}^{p}(\RR f_{*}\mathcal{J}^{n},L) \\
 &= \Hom_{S}^{p}(\RR f_{!}(U,\mathcal{O}_{U}),L).
\end{align*}
La proposition en r\'esulte.

On supposera admis pour $\RR f^{!}$ un formalisme de variance analogue \`a celui du \bookref{app:thm1}{th\'eor\`eme 1} pour $\RR f_{!}$.

Si on veut rendre plus explicite la d\'efinition de $\RR f^{!}$ dans le cas propre, on peut remarquer que si $\mathcal{F}$ quasi-coh\'erent repr\'esente un foncteur $F$, pour tout ouvert $U$ et $\mathcal{J}$ d\'efinissant $X\setminus U$, on a $\mathcal{F}(U) = \varinjlim \Hom(\mathcal{J}^{n}, \mathcal{F}) = \varinjlim F(\mathcal{J}^{n})$.

\section*{n\textdegree{} 5. Calcul de $\RR f^{!}$.}\label{app:sec5}

Pour un morphisme fini, $\RR f^{!}$ est ce qu'on veut; pour le voir, il suffit de prendre pour foncteur r\'esolvant, dans le calcul de $\RR f_{*}$, l'identit\'e. Pour l'espace projectif, l'unicit\'e d'un adjoint ne laisse pas le choix (on utilise le fait qu'on a la formule d'adjonction pour $\RR f_{!} = \RR f_{*}: \DD(X) \to \DD(S)$), on reviendra sur la fl\`eche. Pour une immersion ouverte, $\RR f^{!} = \RR f^{*}$. Pour tout ouvert quasi-projectif $U$ de $X$ au-dessus de $S$, on peut donc calculer la restriction \`a $U$ de $\RR f^{!}$. En particulier, si on peut v\'erifier localement que $\mathcal{H}^{q}\RR f^{!}K = 0$ pour $q \neq p$, on conna\^it $\RR f^{!}K$, puisqu'un objet de la cat\'egorie d\'eriv\'ee n'ayant qu'un faisceau de cohomologie non nul est d\'etermin\'e par ce faisceau. Rappelons un cas important o\`u cette condition est v\'erifi\'ee, avec $K = \mathcal{O}$.

\begin{bookproposition}{7}\label{app:prop7}
Soit $f: X \to Y$, compactifiable. Supposons que localement (sur $X$ et $Y$) on puisse trouver un diagramme
\[
\begin{tikzcd}[column sep=large,row sep=large]
X \ar[r,"f"] \ar[dr,"g"'] & Y \ar[d,"h"] \\
 & S
\end{tikzcd}
\]
o\`u $g$ et $h$ sont localement intersection compl\`ete de dimension relative $d'$ et $d''$, et soit $d = d' - d''$. Alors $\RR f^{!}\mathcal{O}$ est r\'eduit \`a un seul faisceau de cohomologie, situ\'e en degr\'e $-d$, qui est inversible.
\end{bookproposition}

En effet, $\RR g^{!}\mathcal{O} = \RR f^{!}\RR h^{!}\mathcal{O}$, on sait que $\RR g^{!}\mathcal{O}$ et $\RR h^{!}\mathcal{O}$ sont du type indiqu\'e, situ\'es en degr\'e $-d'$ et $-d''$, par des calculs locaux dans l'espace projectif; tout \'etant local, on peut supposer $\mathcal{H}^{-d''}\RR h^{!}\mathcal{O}$ isomorphe \`a $\mathcal{O}$, et l'assertion en r\'esulte.

Pour passer de l\`a \`a des complexes plus g\'en\'eraux, il faudrait au pr\'ealable d\'emontrer des formules du type $\RR f^{!}(K \derivedtensor L) = \RR f^{!}(K) \derivedtensor \LL f^{*}L$. Seule la d\'efinition de la fl\`eche pose un probl\`eme, l'isomorphie \'etant un probl\`eme local. Il faut bien s\^ur des restrictions de degr\'e. Je n'ai rien v\'erifi\'e en d\'etail de ce qui suit.

\subsection*{1\`ere m\'ethode pour d\'efinir la fl\`eche.}

Supposons $K$ et $L$ dans $\DD_{\operatorname{coh}}^{b}(S)$ (et pas seulement quasi-coh\'erents), et en outre que
\begin{enumerate}
\item[(a)] $\RR f^{!}K$ est \`a degr\'e born\'e (il est automatiquement coh\'erent, c'est un probl\`eme local);
\item[(b)] $K \derivedtensor L$ est \`a degr\'e born\'e;
\item[(c)] $\RR f^{!}K \derivedtensor \LL f^{*}L$ est \`a degr\'e born\'e.
\end{enumerate}

Pour d\'efinir
\[
\RR f^{!}K \derivedtensor \LL f^{*}L
 \longrightarrow \RR f^{!}(K \derivedtensor L),
\]
il suffit par dualit\'e de d\'efinir
\[
\RR f_{!}(\RR f^{!}K \derivedtensor \LL f^{*}L)
 \longrightarrow K \derivedtensor L,
\]
et par K\"unneth
\[
\RR f_{!}(\RR f^{!}K \derivedtensor \LL f^{*}L)
 = \RR f_{!}\RR f^{!}K \derivedtensor L;
\]
la fl\`eche cherch\'ee est induite par la fl\`eche d'adjonction
$\RR f_{!}\RR f^{!}K \longrightarrow K$.

\subsection*{2\`eme m\'ethode.}

On suppose $f$ compactifiable, plat et localement d'intersection compl\`ete relative. On a alors $\RR f^{!}\mathcal{O} = \omega[d]$, $\omega$ \'etant inversible. D\'esignons par $\RR f_{!}$ un proc\'ed\'e de calcul de $\RR f_{!}$ au niveau des complexes (c-\`a-d un tel proc\'ed\'e pour un compactifi\'e $\overline{f}: \overline{X} \to S$ de $f$), qu'on suppose du type d\'ecrit au \bookref{app:sec4}{n\textdegree{} 4}, et tel que $C^{p}\mathcal{F}$ soit nul pour $p > d$ (ce qui peut s'obtenir par tronquage). On a alors une fl\`eche
\[
R^{d}\overline{f}_{*}\mathcal{F}[-d]
 \longrightarrow \RR\overline{f}_{*}\mathcal{F}
\]
pour tout $\mathcal{F}$ sur $\overline{X}$ (le dernier faisceau de cohomologie s'envoie dans son complexe). On en tire des morphismes de complexes
\[
\Hom_{S}^{\bullet}(\RR\overline{f}_{*}\mathcal{F}, I)
 \longrightarrow
\Hom_{S}^{\bullet}(R^{d}\overline{f}_{*}\mathcal{F}[-d], I)
\]
pour tout complexe $I$ sur $S$. Soit $I$ un injectif. $R^{d}\overline{f}_{*}\mathcal{F}[-d]$ est exact \`a droite en $\mathcal{F}$, le foncteur en $\mathcal{F}$ \`a droite de la fl\`eche est donc repr\'esentable: on d\'efinit $I_1$ par
\[
\Hom_{S}(R^{d}\overline{f}_{*}\mathcal{F}, I)
 = \Hom_{\overline{X}}(\mathcal{F}, I_1),
 \qquad \RR f_{1}^{!}I = I_1[d].
\]
Ce qui pr\'ec\`ede d\'efinit
\[
\RR f^{!}I \longrightarrow \RR f_{1}^{!}I
\]
et pour un complexe d'injectifs, $\RR f_{1}^{!}$ se d\'efinit terme \`a terme. Il reste \`a \'evaluer $\RR f_{1}^{!}I$ sur des ouverts affines, soit $U$, de $X$; c'est donn\'e par $\varinjlim_{n}\Hom(R^{d}f_{!}\mathcal{J}^{n}, I)$. Ce probl\`eme \'etant local en haut, il est facile de montrer par des m\'ethodes projectives que le dual de
\[
\text{``}\varprojlim_{n}\text{''}\,R^{d}f_{!}\mathcal{J}^{n}
 = R^{d}f_{!}(U,\mathcal{O}),
\]
un Ind-objet, est $f_{*}(U,\omega)$, repr\'esent\'e comme limite de ses sous-modules de sections ayant un p\^ole d'ordre au plus $n$ hors de $U$ ($n \to \infty$).

Ceci donne $\RR f_{1}^{!}I = \omega \otimes f^{*}I[d]$, et la formule g\'en\'erale $\RR f_{1}^{!}K = f^{*}K \otimes \omega[d]$.

\section*{Bibliographie (pour l'Appendice)}
\begin{enumerate}
\item[{[1]}]\phantomsection\label{app:bib1} J.L.\ Verdier, S\'eminaire Bourbaki, Novembre 65, expos\'e 300.
\item[{[2]}]\phantomsection\label{app:bib2} M.\ Nagata, Imbedding of an abstract variety in a complete variety. Journal of Math.\ of Kyoto University, vol.\ 2, n\textdegree{} 1, 1962.
\end{enumerate}

\clearpage
\addcontentsline{toc}{chapter}{Errata}
\chapter*{Errata}
\markboth{Errata}{Errata}

\begin{itemize}
\item Page 10, line 1: read ``see Appendix'' instead of ``unpublished''.
\item Page 85, Proposition 1.1: For the case of $\Qco(X)$, one needs the hypothesis ``$X$ quasi-separated''.
\item Page 88, line 5: read ``proper'' instead of ``of finite type''.
\item Page 120, Ch II \S7: It has recently come to my attention that S.\ Kleiman has independently arrived at the results of [II 7.8] and [II 7.11] (unpublished).
\item Page 137, line 5: refer to [EGA IV \S16].
\item Page 144, line 4: read ``invertible sheaf shifted $n$ places to the left, where $n = \operatorname{rel.\ dim} X/Y$'' instead of ``invertible sheaf''.
\item Page 185, line 9: observe that the condition ``$f$ finite'' is a consequence of the other conditions.
\item Page 190, line 7: read ``a unique isomorphism'' instead of ``an isomorphism''.
\item Page 199: add at bottom: ``Furthermore, the residue symbol is uniquely characterized by the properties (R0), (R1), (R2), (R5), (R6), and (R7).''
\item Page 249, line 8: ``is'' instead of ``in''.
\item Page 254, line 8: ``$\mathbb{Z}$'' instead of ``$\mathbb{Q}$''.
\item Page 288, Theorem 8.3: One must assume also that the fibres of $f$ are of bounded dimension, so that $f^{\sharp}(R^{\bullet})$ will be in $\DD^{+}$.
\item Page 297, line 7: read ``strengthens'' instead of ``generalizes''.
\item Page 298, line 7: insert ``shifted'' after ``invertible sheaf''.
\item Page 301, line 5: ``not'' instead of ``now''.
\item Page 301, Problem 2: Giraud [6] provides us with an element of
\[
H^{2}(X, \mathbb{G}_{m})
\]
(cohomology in the Zariski topology), whose vanishing is a necessary and sufficient condition for the global existence of a dualizing complex, supposed to exist locally. On the other hand, it is easy enough to construct a noetherian scheme of finite Krull dimension, which has a dualizing complex locally, but none globally, due to the lack of a global codimension function.
\item Page 373, line 1: change reference to [EGA V].
\item Page 403, last line: read ``Publ.\ Math.\ I.H.E.S.\ n\textdegree{} 29''.
\end{itemize}